% process_flow.tex - Process flow diagram macros for Paperforge
% Requires: _colors.tex, _styles.tex loaded first
%
% Macros provided:
%   \processflow{step1}{step2}{step3}           - 3-step horizontal flow
%   \processflowfour{s1}{s2}{s3}{s4}            - 4-step horizontal flow
%   \processflowvertical{step1}{step2}{step3}   - 3-step vertical flow
%   \processflowlabeled{s1}{l1}{s2}{l2}{s3}     - Flow with connector labels
%   \processflowdark{step1}{step2}{step3}       - Dark theme 3-step flow
%
% Cognitive load: Maximum 4 steps per diagram (RAG principle)

% ============================================================================
% THREE-STEP HORIZONTAL FLOW
% ============================================================================
% Usage: \processflow{Input}{Process}{Output}
\newcommand{\processflow}[3]{%
\begin{center}
\begin{tikzpicture}[node distance=3.5cm]
    \node[diagram box] (a) {#1};
    \node[diagram box, right=of a] (b) {#2};
    \node[diagram box, right=of b] (c) {#3};
    \draw[diagram arrow] (a) -- (b);
    \draw[diagram arrow] (b) -- (c);
\end{tikzpicture}
\end{center}
}

% ============================================================================
% FOUR-STEP HORIZONTAL FLOW
% ============================================================================
% Usage: \processflowfour{Step 1}{Step 2}{Step 3}{Step 4}
\newcommand{\processflowfour}[4]{%
\begin{center}
\begin{tikzpicture}[node distance=2.8cm]
    \node[diagram box, minimum width=2cm] (a) {#1};
    \node[diagram box, minimum width=2cm, right=of a] (b) {#2};
    \node[diagram box, minimum width=2cm, right=of b] (c) {#3};
    \node[diagram box, minimum width=2cm, right=of c] (d) {#4};
    \draw[diagram arrow] (a) -- (b);
    \draw[diagram arrow] (b) -- (c);
    \draw[diagram arrow] (c) -- (d);
\end{tikzpicture}
\end{center}
}

% ============================================================================
% THREE-STEP VERTICAL FLOW
% ============================================================================
% Usage: \processflowvertical{Top}{Middle}{Bottom}
\newcommand{\processflowvertical}[3]{%
\begin{center}
\begin{tikzpicture}[node distance=1.5cm]
    \node[diagram box] (a) {#1};
    \node[diagram box, below=of a] (b) {#2};
    \node[diagram box, below=of b] (c) {#3};
    \draw[diagram arrow] (a) -- (b);
    \draw[diagram arrow] (b) -- (c);
\end{tikzpicture}
\end{center}
}

% ============================================================================
% LABELED FLOW (with connector labels)
% ============================================================================
% Usage: \processflowlabeled{Step1}{label1}{Step2}{label2}{Step3}
\newcommand{\processflowlabeled}[5]{%
\begin{center}
\begin{tikzpicture}[node distance=3.5cm]
    \node[diagram box] (a) {#1};
    \node[diagram box, right=of a] (b) {#3};
    \node[diagram box, right=of b] (c) {#5};
    \draw[diagram arrow] (a) -- node[diagram label above] {#2} (b);
    \draw[diagram arrow] (b) -- node[diagram label above] {#4} (c);
\end{tikzpicture}
\end{center}
}

% ============================================================================
% DARK THEME THREE-STEP FLOW
% ============================================================================
% Usage: \processflowdark{Step1}{Step2}{Step3}
\newcommand{\processflowdark}[3]{%
\begin{center}
\begin{tikzpicture}[node distance=3.5cm]
    \node[diagram box dark] (a) {#1};
    \node[diagram box dark, right=of a] (b) {#2};
    \node[diagram box dark, right=of b] (c) {#3};
    \draw[diagram arrow dark] (a) -- (b);
    \draw[diagram arrow dark] (b) -- (c);
\end{tikzpicture}
\end{center}
}

% ============================================================================
% DARK THEME LABELED FLOW
% ============================================================================
% Usage: \processflowdarklabeled{Step1}{label1}{Step2}{label2}{Step3}
\newcommand{\processflowdarklabeled}[5]{%
\begin{center}
\begin{tikzpicture}[node distance=3.5cm]
    \node[diagram box dark] (a) {#1};
    \node[diagram box dark, right=of a] (b) {#3};
    \node[diagram box dark, right=of b] (c) {#5};
    \draw[diagram arrow dark] (a) -- node[diagram label dark, above, midway] {#2} (b);
    \draw[diagram arrow dark] (b) -- node[diagram label dark, above, midway] {#4} (c);
\end{tikzpicture}
\end{center}
}

% ============================================================================
% CYCLIC FLOW (3 steps in a loop)
% ============================================================================
% Usage: \processcycle{Step1}{Step2}{Step3}
\newcommand{\processcycle}[3]{%
\begin{center}
\begin{tikzpicture}
    \node[diagram box] (a) at (90:2cm) {#1};
    \node[diagram box] (b) at (210:2cm) {#2};
    \node[diagram box] (c) at (330:2cm) {#3};
    \draw[diagram arrow] (a) to[bend right=15] (b);
    \draw[diagram arrow] (b) to[bend right=15] (c);
    \draw[diagram arrow] (c) to[bend right=15] (a);
\end{tikzpicture}
\end{center}
}

% ============================================================================
% BRANCHING FLOW (1 to 2)
% ============================================================================
% Usage: \processbranch{Start}{Branch A}{Branch B}
\newcommand{\processbranch}[3]{%
\begin{center}
\begin{tikzpicture}[node distance=2cm]
    \node[diagram box] (start) {#1};
    \node[diagram box, below right=1.5cm and 2cm of start] (a) {#2};
    \node[diagram box, above right=1.5cm and 2cm of start] (b) {#3};
    \draw[diagram arrow] (start) -- (a);
    \draw[diagram arrow] (start) -- (b);
\end{tikzpicture}
\end{center}
}

% ============================================================================
% MERGING FLOW (2 to 1)
% ============================================================================
% Usage: \processmerge{Input A}{Input B}{Output}
\newcommand{\processmerge}[3]{%
\begin{center}
\begin{tikzpicture}[node distance=2cm]
    \node[diagram box] (a) at (0, 1.2) {#1};
    \node[diagram box] (b) at (0, -1.2) {#2};
    \node[diagram box] (out) at (4.5, 0) {#3};
    \draw[diagram arrow] (a) -- (out);
    \draw[diagram arrow] (b) -- (out);
\end{tikzpicture}
\end{center}
}

% ============================================================================
% FEEDBACK LOOP
% ============================================================================
% Usage: \processfeedback{Process}{Feedback}{Output}
\newcommand{\processfeedback}[3]{%
\begin{center}
\begin{tikzpicture}[node distance=3cm]
    \node[diagram box] (process) {#1};
    \node[diagram box, right=of process] (output) {#3};
    \node[diagram box, below=1.5cm of process] (feedback) {#2};
    \draw[diagram arrow] (process) -- (output);
    \draw[diagram arrow] (output) |- (feedback);
    \draw[diagram arrow] (feedback) -- (process);
\end{tikzpicture}
\end{center}
}
